Sets are the backbone of mathematics and we have been using sets throughout this course.  In this chapter, we introduce formal set theory.  We begin with set vocabulary and operations in \Cref{sec_sets} {\secsets}.  We revisit subsets in \Cref{sec_subsets} {\secsubsets} and work with sets whose elements are sets. We also look at set-theoretic functions.  In \Cref{sec_partitions} {\secpartitions}, we introduce partitions of sets and numbers and introduce a dynamical system on a state graph related to partitions.  In \Cref{sec_relations} {\secrel} we define the Cartesian product of sets and look at functions involving Cartesian products of sets.  We also introduce relations, an abstraction that generalizes functions, operations, and relationships such as inequalities.  Last, we explore properties of relations, equivalence relations, and equivalence classes in \Cref{sec_equiv_relations} {\secequivrel}.